\documentclass[11pt,a4paper]{article}
\usepackage[T1]{fontenc}
\usepackage[utf8]{inputenc}
\usepackage{lmodern}
\usepackage[a4paper,margin=27mm]{geometry}
\usepackage{amsmath,amssymb,amsthm,mathtools}
\usepackage{needspace}
\usepackage[expansion=false]{microtype}
\usepackage{xurl}
\usepackage[colorlinks=true,linkcolor=blue,citecolor=blue,urlcolor=blue]{hyperref}
\allowdisplaybreaks[2]
\setlength{\emergencystretch}{3em}
\numberwithin{equation}{section}
\newtheorem{theorem}{Theorem}[section]
\newtheorem{proposition}[theorem]{Proposition}
\newtheorem{lemma}[theorem]{Lemma}
\theoremstyle{remark}
\newtheorem{remark}[theorem]{Remark}
\newcommand{\Rea}{\operatorname{Re}}
\newcommand{\src}[1]{\texttt{\detokenize{#1}}}
\hypersetup{
  pdftitle={A Parameter Refinement of the Seven-Eighths Zero-Free Half-Plane},
  pdfauthor={GPT-6 Astra, Joseph M. Shunia}
}
\title{A Parameter Refinement of the Seven-Eighths\\Zero-Free Half-Plane}
\author{GPT-6 Astra, Joseph M. Shunia}
\date{7 October 2026}

\begin{document}
\maketitle
\begin{abstract}
We refine the parameter choices in OpenAI's claimed zero-free half-plane
$\Rea s>7/8$ for finite-order Hecke $L$-functions over
$\mathbb Q(\sqrt{-3})$. Here $s$ is a complex number and $\Rea s$ is its
real part, so the claimed region consists of those $s$ whose real part
is greater than $7/8$. Taking the source's arithmetic and analytic
estimates as inputs, we obtain the sharper boundary
\[
 \sigma_*=0.874957019420098946\ldots,
 \qquad
 7884\sigma_*^3-18819\sigma_*^2+14643\sigma_*-3686=0.
\]
Since $\sigma_*<162861/186136=7/8-1/23267$, the half-plane
$\Rea s>162861/186136$ is zero-free as well.
The argument retains the exact zero-free-parameter dependence in a
selected-prime fourth-moment bound, extends that bound to the required
parameter range, and optimizes the geometry of the existing compensated
probe. We verify the resulting endpoint inequality and the affected
Euler-product and contour estimates, obtaining a contracting bootstrap
map with fixed point $\sigma_*$. The conclusion transfers, as in the
source, to all Dirichlet $L$-functions. This is a \emph{source-relative}
refinement: it neither reproves nor independently validates the deep
estimates in the cited manuscript.
\end{abstract}

\noindent\textbf{Keywords.} Zero-free regions; Dirichlet $L$-functions;
Hecke $L$-functions; compensated probes; moment estimates.

\section{Introduction and main result}
\label{sec:statement}

The compensated-probe method of \cite{OA} is reported to give the
uniform zero-free region $\Rea s>7/8$ for finite-order Hecke
$L$-functions over $\mathbb Q(\sqrt{-3})$, with a transfer to Dirichlet
$L$-functions. Its main input is a pair of estimates for a finite
arithmetic probe: a direct bound and a Mellin representation whose
principal signal contains a reciprocal $L$-function. A zero to the
right of the permitted boundary would contradict the continuation
obtained from these estimates.

The purpose of this note is narrower than establishing those inputs.
We preserve the source's finite probe and its deep reflection, detector,
and moment identities, but sharpen the treatment of the selected-prime
capacity and optimize the allowed asymmetric scales. A parameter-dependent
fourth-moment estimate, an exact endpoint inequality, and an audit of the
affected Euler and contour regions make the resulting improvement
iterable. The gain in the boundary is small, but strict.


Let $F=\mathbb Q(\sqrt{-3})$. Write $\beta_*$ for the supremum of
$1/2$ and of the real parts of zeros in $1/2\le\Rea s\le1$ of all
primitive finite-order Hecke $L$-functions over $F$. Poles are excluded.
The source \cite{OA} proves $\beta_*\le7/8$. Its transfer proposition
passes a bound on this family to all finite-order Hecke functions and
all Dirichlet $L$-functions, including $\zeta$.

\Needspace{10\baselineskip}
\begin{theorem}[Source-relative refinement]\label{thm:main-refinement}
Assume the results of \cite{OA} listed below, with their stated
uniformity and coefficient conventions. Let $\sigma_*$ be the unique
root in $(17499/20000,7/8)$ of
\begin{equation}\label{eq:main-cubic}
 7884s^3-18819s^2+14643s-3686=0.
\end{equation}
Then $\beta_*\le\sigma_*$. In particular every finite-order Hecke
$L$-function over $F$, and every Dirichlet $L$-function, is zero-free
in the strict half-plane $\Rea s>\sigma_*$. A principal pole at one
is allowed. No assertion of nonvanishing on $\Rea s=\sigma_*$ is made.
Since
\begin{equation}\label{eq:main-rational}
 \sigma_*<\frac{162861}{186136}=\frac78-\frac1{23267},
\end{equation}
these functions are, in particular, zero-free in
$\Rea s>162861/186136$.
\end{theorem}

\paragraph{Dependence on the source.}
We use the $7/8$ theorem and the exact compensated probe and Poisson
identity from \cite{OA}; its cubic-theta reflection, finite-correlation
and additive Gram estimates; its buffered zero detector, strict marked
inverse moment and sixth-power amplification; and the parameter-independent
parts of its plain fourth-moment argument. We also invoke the source's
logarithmic control, fixed-ray prime normalization, contour accounting,
principal-residue, outer-tail, late-height and transfer results. The
changes to their relevant parameter ranges are justified in the appendices.
In particular, no uniform theorem for arbitrary moving coefficients or
punctures is taken as an additional input.

Source equation and lemma labels printed in typewriter font refer to the
frozen version of \cite{OA} identified in the bibliography.

\section{The improvement map and compensated probe}
\label{sec:map}
\subsection{Parameters and bootstrap}

For an already known upper bound
\[
 \beta_*\le\sigma,\qquad
 \frac{17499}{20000}\le\sigma\le\frac78,
 \qquad \kappa=2\sigma-1,
\]
define
\begin{align}
 d_\kappa&=
 \frac{156\kappa-19-
 \sqrt{9936\kappa^2-2808\kappa+201}}
 {12(12\kappa-1)},\label{eq:main-critical-bin}\\
 \ell&=\frac{5-6d_\kappa}{9+18d_\kappa},\label{eq:main-ell}\\
 b&=(1+3\ell)
 \frac{36\kappa d_\kappa-6\kappa-1}
 {(156\kappa-19)-12(12\kappa-1)d_\kappa},\label{eq:main-b}\\
 l_x&=\frac{1-\ell-b}{2},\qquad
 l_y=\frac{1-\ell+b}{2},\qquad
 h=\frac{1+3\ell+b}{2}.\label{eq:main-geometry}
\end{align}
Thus $l_x+l_y+\ell=1$ and $h=1-l_x+\ell$.
The symbol $d_\kappa$ is one critical zero-bin parameter; it is not
the row-scale exponent $d$ used in the estimates below.
The improvement map is
\begin{equation}\label{eq:main-map}
 f(\sigma)=\frac{11}{12}-\frac\ell4
          =1-\frac{2}{9+18d_{2\sigma-1}}.
\end{equation}
\begin{proposition}[Bootstrap step]\label{prop:bootstrap}
Under the source hypotheses of Theorem~\ref{thm:main-refinement},
for every $\sigma\in[17499/20000,7/8]$,
\begin{equation}\label{eq:main-bootstrap}
 \beta_*\le\sigma\quad\Longrightarrow\quad
 \beta_*\le f(\sigma).
\end{equation}
\end{proposition}
The fixed-point calculation in Appendix~\ref{sec:algebra-audit}
shows that $f$ has a unique fixed point $\sigma_*$ in this interval
and satisfies $f(\sigma)<\sigma$ above that point. Thus
Proposition~\ref{prop:bootstrap}, iterated from $7/8$, implies
Theorem~\ref{thm:main-refinement}.

\subsection{The unchanged arithmetic probe}

For clarity, we retain the source's finite arithmetic observable rather
than defining a new probe by its desired Mellin representation. Put
$X=Z^{l_x}$, $Y=Z^{l_y}$, and choose disjoint fixed-ray prime windows
of lengths $P_i=Z^{\ell_i}$, where $\sum_i\ell_i=\ell$.
Write $I_{\eta;D}$ for the source's completed cubic-theta probe with
$1_{D\mid cn^3}$ inserted. For $p_J=\prod_{i\in J}p_i$, the physical
probe is exactly
\begin{align}
 I_{\eta,\mathrm{modified}}(Z)
 ={}&\sum_{(p_i)}\prod_i W_i(q_{p_i}/P_i)
 \sum_{J\subseteq\{1,\ldots,K\}}
 (-1)^{|J|}q_{p_J}^{-3/2}\overline{\eta(p_{J^c})}\notag\\[-1mm]
 &\hspace{22mm}\cdot
 I_{\eta;p_{J^c}}(X/q_{p_J},Y/q_{p_J},Zq_{p_{J^c}}).
 \label{eq:physical-probe}
\end{align}
Here $q_a=Na$, every slot retains its original window under all
rescalings, and all source zero extensions and fixed-ray phases remain.
Only its real exponents and the number of fixed slots are changed.
This finite expression is the source's
\src{eq:compensated-probe-definition}.

At a selected prime its two terms cancel the scalar local contribution,
leaving the character prime factor to which the inverse and plain
moments apply. The gain proved here comes from keeping the exact
prime capacity and changing the geometry around that cancellation.
No new off-diagonal cancellation estimate is assumed.




\subsection{Proof structure}
The refinement has three inputs beyond the assertions imported from
\cite{OA}. Appendix~\ref{sec:extended-moment-audit} extends the plain
fourth moment and derives a selected-row count without the source's
fixed-capacity comparison loss. Appendix~\ref{sec:algebra-audit}
optimizes the geometry and proves the global endpoint inequality by
rational coefficient bounds; its fixed-point calculation identifies
$\sigma_*$. Appendix~\ref{sec:analytic-audit} verifies the perturbed
low estimate, Euler-product region, contour inequalities, and Mellin
continuation with the needed quantifier order. The bootstrap is
assembled in Section~\ref{sec:assembly}.

\section{Proof of the refinement}
\label{sec:assembly}

\begin{proof}[Proof of Proposition~\ref{prop:bootstrap} and Theorem~\ref{thm:main-refinement}]
Fix an input $\sigma\in[17499/20000,7/8]$ with $\beta_*\le\sigma$,
and take \eqref{eq:main-critical-bin}--\eqref{eq:main-geometry}.
The estimates in Appendices~\ref{sec:algebra-audit} and
\ref{sec:analytic-audit} establish the needed geometric and Euler
inequalities. Suppose, for a
contradiction, that
\[
 \Delta=\beta_*-f(\sigma)>0.
\]
This is a fixed positive number for the entire Hecke family. It need
not be effective or known numerically, and it does not depend on the
target character.

The extended fourth moment from Appendix~\ref{sec:extended-moment-audit}
applies because
$\beta_*\le(1+\kappa)/2$. The buffered dynamic bins satisfy
$\delta\le2\beta_*-1\le\kappa$. The exact selected-row bound therefore
applies with the same $\kappa$ and gives the envelope
\eqref{eq:audit-Rstar}; it is not permissible to insert the desired
new bound in place of this already available input. The global
algebraic inequality gives $E(h)\le0$. The floor, outer, and principal
contributions have the explicit strict margins checked above.

Choose the positive moment losses, bin and amplitude widths, capacity
decrements, and row extension so that their combined real cost is
smaller than a fixed fraction of $\Delta$ and of the other strict
margins. Choose a sufficiently fine fixed prime mesh and an even number
of disjoint slot windows. Their existence uses the positive supply
margin, and their normalizer is nonzero for all sufficiently large
$Z$. Every one of these real choices precedes the target character.
The source's coefficient classes remain intact because each selected
prime coefficient is a fixed finite sum of fixed-ray characters,
independent of the varying row.

For the chosen target $\eta$, fix its arithmetic data and excluded set.
Fix the finite internal seminorm and height orders first, then choose
the height exponent small enough to absorb them, and only afterwards
choose the external tail order. The late-height argument in the
analytic audit supplies a positive saving independent of $\eta$.
With the same scalar normalizer on both sides, it follows that
\[
 |J_\eta(Z)|\ll_\eta Z^{C(f(\sigma))+\Delta/2},\qquad
 |J_\eta(Z)-f_\eta(Z)|
       \ll_\eta Z^{C(\beta_*)-m}
\]
for one $m>0$ independent of $\eta$. The notation $f_\eta$ denotes the
Mellin signal, whereas $f(\sigma)$ denotes the numerical improvement
map. These are different objects.

The Euler multiplier of the signal is holomorphic and bounded away
from zero on $\Rea s>f(\sigma)$. Its Mellin transform consequently
continues every reciprocal $1/L_F^S(s,\eta)$ into a common half-plane
strictly to the left of $\beta_*$. The detailed justification in the
analytic audit does not move a contour through a presumed zero. This
contradicts the definition of $\beta_*$ and proves
\eqref{eq:main-bootstrap}.

Starting with the source theorem $\beta_*\le7/8$, successive
applications give $\beta_*\le f^{\circ n}(7/8)$ for every finite $n$.
The algebraic contraction shows that these numbers decrease to
$\sigma_*$. Hence $\beta_*\le\sigma_*$. No common analytic constant
across infinitely many iterations is required. Equivalently,
if $\beta_*>\sigma_*$, apply the same proof directly with
$\sigma=\beta_*$, since $f(\beta_*)<\beta_*$.
The source transfer proposition now gives the Dirichlet conclusion
in Theorem~\ref{thm:main-refinement}. This completes the source-relative proof.
\end{proof}

\section{Scope of the refinement}
\label{sec:scope}
The exact $\kappa$-dependent selected-prime capacity, the perturbed
low estimate, the extended local Euler region, and the global endpoint
inequality are all required for Proposition~\ref{prop:bootstrap}.
Optimizing the numerical endpoint alone would not yield a zero-free
improvement: in particular, the positive-part error for rescaled
subsets in the low row estimate must be accounted for before summing
the selected-prime tuples.

The exponent-ledger comparison in Appendix~\ref{sec:algebra-audit}
is an optimality statement only within the particular moment bounds
and probe family used here. It does not preclude stronger zero-free
regions from new arithmetic estimates, a different compensated probe,
or a different treatment of the high-frequency terms. None of the
deep estimates imported from \cite{OA} is independently established
in this note.

\appendix

\section{Fourth-moment extension and selected-row estimate}
\label{sec:extended-moment-audit}

The assertions in this section are extensions of specified results of the
source manuscript, not new proofs of its finite-correlation identities or
its unmarked inverse moment.  We use those source results as mathematical
inputs.  The dependence on the changed parameter is checked below by explicit
inequalities; no assertion follows merely by continuity in that parameter.

\subsection{Source hypotheses and notation}
Retain the arithmetic data, row families, natural zero extensions, fixed
finite group $\Theta$, plain polynomials $S_{\psi_k}(n;W)$, and selected
prime polynomials of source Lemma \texttt{lem:plain}.  Write $M$ for its
effective row width, $z$ for the total length of the selected prime factors,
and $A=n_1+n_2+z$.  Positive prime-factor coefficients have the fixed
finite-ray expansion required in that lemma; no row-dependent coefficient
is admitted.  Besides the parameter-independent calculations in the proof
of \texttt{lem:plain}, we use the global logarithmic-derivative estimate of
\texttt{lem:logarithmic-control}, the strict marked inverse moment
\texttt{lem:marked}, sixth-power amplification
\texttt{lem:inverse-amplification}, and the two witnesses in
\texttt{prop:detector-witness}.

\begin{lemma}[Extended plain fourth moment]\label{lem:plain-extension}
Source Lemma \texttt{lem:plain}, with all its coefficient, support,
uniformity, and row-family hypotheses, remains valid for
\begin{equation}\label{eq:audit-kappa-range}
             \frac{13}{18}\le\kappa\le1.
\end{equation}
More explicitly, its conclusion is
\begin{equation}\label{eq:audit-plain-moment}
 \sum_{k\in\mathcal R_z}
 |S_{\psi_k}(n_1;W_1)S_{\psi_k}(n_2;W_2)Q|^2
 \ll Z^{M+\varepsilon}.
\end{equation}
For $z=0$ its plain lengths are unrestricted in the prescribed bounded
ranges.  For $z>0$ assume
\begin{equation}\label{eq:audit-affine-condition}
 A+(6\kappa-1)z=n_1+n_2+6\kappa z\le M,
\end{equation}
and, when $\kappa<1$, assume $\beta_*\le(1+\kappa)/2$.
The slot mesh, exponent losses, and finite seminorm/height assertions can
be chosen uniformly for \eqref{eq:audit-kappa-range}.

\end{lemma}
\begin{proof}
There are four substantive uses of $\kappa$ in the source proof and one
uniform Lipschitz bound.  We verify each, retaining the source induction
and its numerical loss budget.

(i) Put $s_\kappa=(1+\kappa)/2$.  The global assertion of source
\texttt{lem:logarithmic-control} applies for every fixed $b\ge\beta_*$,
not only for $b\ge7/8$.  Thus, on $\Re s=s_\kappa+e$,
\[
 \frac{L'}{L}(s,\psi)\ll_e
 \log\{2Q_\psi(3+|\Im s|)^2\}
\]
for each nonprincipal primitive inducing character.  The exclusion
of the inducing characters in $\Theta$ ensures that twisting by each
fixed prime-factor coefficient character is still nonprincipal.  Mellin
inversion and contour shifting give the source prime-factor estimate
\begin{equation}\label{eq:audit-prime-cost}
 |Q_{\boldsymbol\omega}|^2
 \ll Z^{\kappa z+\varepsilon_1}
 p_j(\boldsymbol W)^b(1+|\boldsymbol\omega|)^h.
\end{equation}
The exponent is $2s_\kappa-1=\kappa$.  Prime powers contribute only
$P^{1/2+\varepsilon_1}$ before central normalization.  Contour displacement
costs $2ez$, and $z$ stays in a bounded range.  On the new compact range,
$s_\kappa\ge31/36>1/2$, so the source proof supplies the same uniform
choice of $e$, losses, and fixed height orders.  For $\kappa=1$ use
absolute prime counting, exactly as in the source.

(ii) At the width floor, $A\ge z$ and
\eqref{eq:audit-affine-condition} imply
\begin{equation}\label{eq:audit-terminal-bound}
 z\le\frac{M}{6\kappa}\le\frac{3M}{13},
             \qquad \kappa z\le\frac M6.
\end{equation}
The source displays $z\le2M/9$ at this point, but its terminal exponent
uses only the second inequality.  Therefore its additional terminal
loss remains $\rho/6$ for a width floor $M\le\rho$.

(iii) In the centered comparison, $A>5M/6$ and
\eqref{eq:audit-affine-condition} imply
\[
 (6\kappa-1)z\le M-A<M/6.
\]
Since $6\kappa-1\ge10/3$, this gives the strict bound
$z<M/20$.  This replaces the source's intermediate inequality
$z<M/21<M/20$ and is exactly the bound used afterwards.  Its reflected
comparison therefore still satisfies
\begin{align}
 A_{\rm comp}&\le\frac32M-A+2z+\xi
                         \le\frac{23}{30}M+\xi,\label{eq:audit-comparison-a}\\
 A_{\rm comp}+(6\kappa-1)z
 &\le\frac52M-2A+2z+\xi
                         \le\frac{14}{15}M+\xi.
                         \label{eq:audit-comparison-b}
\end{align}
Also $A-z>5M/6-M/20>M/2$, as required for the two comparison lengths.
Both relevant boundaries retain the source margin $M/15$ before $\xi$;
its choice $\xi\le\rho/30$ therefore works unchanged.

(iv) At a smaller-width child define its positive affine defect by
\[
 F_{\rm act}=
 (A_{\rm act}-M'_{\rm act}+(6\kappa-1)z_{\rm act})_+.
\]
Removing a whole prime factor of length $d$ decreases this affine
quantity by $6\kappa d$ and, by \eqref{eq:audit-prime-cost}, costs
$\kappa d$ in the squared exponent.  The first prefix reaching
$F_{\rm act}/(6\kappa)$ overshoots by at most the mesh $\eta$.
If no such prefix exists, all slots are removed and the unrestricted
zero-slot child is used.  Consequently
\begin{equation}\label{eq:audit-greedy}
 \kappa d_z\le F_{\rm act}/6+\kappa\eta
                         \le F_{\rm act}/6+\eta.
\end{equation}
This is the source's exact charge, with the same numerical constant
$1/6$ and the same single overshoot.  No lower bound $\kappa\ge3/4$
is used here.

(v) The source aggregate boundary-error estimate uses only
$0\le6\kappa-1\le5$, through
\[
 |\Delta M|+|\Delta A|+5|\Delta z|.
\]
It remains valid.  All other transformations in the source proof are
independent of $\kappa$: fixed-mask removal, reflection of at most two
plain factors, finite Gauss identities, the two Poisson steps, the
exceptional-character count, the equal-product-scale centered
coefficient identity, and the width decrease.  In particular the new
range changes neither the depth bound nor the source terminal budget
$\rho+\delta+\rho/6+3\sigma+5\sigma/3$.
The same choices of $\xi,\eta,\varepsilon_0$ then close its finite
induction.  This proves \eqref{eq:audit-plain-moment} under the stated
source inputs.

\end{proof}
\subsection{The selected-row bound}
For the rest of this section assume
\begin{equation}\label{eq:audit-count-range}
 \frac{13}{18}\le\kappa\le\frac34,
 \qquad \beta_*\le\frac{1+\kappa}{2},
 \qquad \alpha=\frac56.
\end{equation}
This includes the narrower interval $[7499/10000,3/4]$ used in the
bootstrap.  Consider one retained dynamic and amplitude bin of
sixth-power-free physical rows $q_u\asymp U$.  Require the source's
actual witnesses, so $a>51/100$, and put
\[
 \delta=2a-1\le\kappa,\qquad
 0\le q\le\delta/2,\qquad x=q/\delta,\qquad
 c=\frac1{3\kappa}.
\]
The equality $\kappa=2\beta_*-1$ is unnecessary: the displayed upper
bound on $\beta_*$ gives both the positive-slot moment and
$\delta\le\kappa$.  Define
\begin{equation}\label{eq:audit-DP}
 D_x=3-(1+2c)x,\qquad P_x=2(1-cx)(1-x).
\end{equation}
If the available prime length exceeds $7/37$ by a fixed positive
amount, then for every $1\le t\le3/2$ the source selection mechanism gives
\begin{equation}\label{eq:audit-row-count}
 \#\mathcal B\ll
 U^{\max\{R_{\rm short}(t),L(t)\}+\varepsilon}
 (1+T_1)^{A_{\mathcal A}},
\end{equation}
where
\begin{align}
 R_{\rm short}(t)&=1-\delta+
       \frac{\delta P_x}{D_x}(3/2-t),\label{eq:audit-short}\\
 L(t)&=1-\delta+(\alpha-\delta)(t-1).
 \label{eq:audit-long}
\end{align}
There is no additional $\Delta/4$ loss.

\paragraph{Proof of the selected-row bound.}
The source witness supplies $r+m\ge t-O(\varepsilon)$ and the individual
spikes $|M_r|^2\gg U^{\delta r-\varepsilon}$,
$|S_m|^2\gg U^{\delta m-\varepsilon}$.  Slot selection and the moment
estimates give respectively
\[
 U^{1-\delta r-2qz+\varepsilon},\qquad
 U^{1-2\delta m-2qz+\varepsilon}
\]
as row-count bounds, with the common height factor understood.  The second
uses two copies of the plain witness.  As in the source, finitely many
presentations and dyadic pairs are separated first; rowwise smooth
parameters are handled by the common Sobolev argument.  These steps and
the coefficient-class verification do not depend on $\kappa$.

For $r<1$ and $m<1/2$, use the exact capacities
\[
 z_M(r)=(1-r)/2,\qquad
 z_P(m)=\frac{1-2m}{6\kappa}=\frac c2(1-2m).
\]
The ideal comparison lines are therefore
\begin{align}
 A_I(r)&=1-\delta\{x+(1-x)r\},\label{eq:audit-AI}\\
 S_t(r)&=1-\delta\{cx+2(1-cx)(t-r)\}.
 \label{eq:audit-SP}
\end{align}
The actual plain exponent decreases with $m$, because its derivative is
$-2\delta(1-cx)<0$.  Thus replacing $m$ by $t-r-O(\varepsilon)$
is legitimate.  The two lines cross at
\begin{equation}\label{eq:audit-crossing}
 r_*(t)=\frac{2(1-cx)t-(1-c)x}{D_x}.
\end{equation}
Here $4/9\le c\le6/13$, and $D_x\ge53/26>0$.
At the endpoints,
\[
 r_*(3/2)=1,\qquad
 r_*(1)=\frac{2-(1+c)x}{3-(1+2c)x}.
\]
The derivative of the last quotient with respect to $x$ has numerator
$c-1<0$, so
\[
 r_*(t)\ge r_*(1)\ge
 \frac{3-c}{5-2c}
 =\frac{9\kappa-1}{15\kappa-2}\ge\frac{23}{37}.
\]
For the final inequality, the derivative of
$(9\kappa-1)/(15\kappa-2)$ is
$-3/(15\kappa-2)^2$, and $\kappa\le3/4$.
Also
\[
 t-r_*(t)=\frac{(1-x)t+(1-c)x}{D_x},\qquad
 \frac13\le t-r_*(t)\le\frac12.
\]
Indeed this expression increases with $t$, equals $1/2$ at $t=3/2$,
and its value at $t=1$ minus $1/3$ is
$(1-c)x/(3D_x)\ge0$.

Choose the plain count for $r\le r_*(t)$ and the inverse count for
$r\ge r_*(t)$.  On the inverse side, requesting
$z=z_M(r)-\nu_0$ gives
\[
 1-r-2z=2\nu_0,\qquad
 3-2r-8z=2r-1+8\nu_0\ge9/37+8\nu_0.
\]
These are the two strict marked-moment margins.  Its requested capacity
is at most $7/37$.  On the plain side $m\ge1/3-O(\varepsilon)$, so
\[
 z_P(m)\le\frac1{18\kappa}+O(\varepsilon)
                         \le\frac1{13}+O(\varepsilon)<\frac7{37}
\]
after reducing preliminary losses.  A fixed supply margin therefore
suffices uniformly.  Decreasing capacities by $\nu_0$, selecting whole
slots from a sufficiently fine mesh, and sorting positive amplitude
slots exactly as in the source cost only $O(\varepsilon)$.

The weighted average
\[
 \frac{2(1-cx)A_I(r)+(1-x)S_t(r)}{D_x}
\]
is independent of $r$ and equals \eqref{eq:audit-short}.
Since $A_I$ decreases and $S_t$ increases, the selected short count is
bounded by this common crossing value.  No approximation to
$1/(6\kappa)$ has been made, so the source comparison loss
\texttt{eq:capacity-comparison} is absent.

For $r\ge1$, sixth-power amplification gives
$1-\alpha+(\alpha-\delta)r\le L(t)+O(\varepsilon)$ since
$\delta\le\kappa<\alpha$ and $r\le t+O(\varepsilon)$.
For $m\ge1/2$, the unweighted plain moment gives $1-\delta$.
When a capacity is at most $\nu_0$, the corresponding unweighted moment
gives $1-\delta+O(\nu_0)$: explicitly $2\delta\nu_0$ on the inverse
side and at most $6\delta\nu_0$ on the plain side.  This covers zero
capacities without invoking a marked moment with a shrinking strict
margin.  All these bounds prove \eqref{eq:audit-row-count}.

Balancing \eqref{eq:audit-short} and \eqref{eq:audit-long} gives
\begin{equation}\label{eq:audit-optimal-t}
 t=1+\frac{\delta P_x}
          {2\{(\alpha-\delta)D_x+\delta P_x\}}\in[1,3/2]
\end{equation}
and therefore the exact envelope
\begin{equation}\label{eq:audit-Rstar}
 R_\kappa(\delta,x)=1-\delta+
 \frac{(\alpha-\delta)\delta P_x}
      {2\{(\alpha-\delta)D_x+\delta P_x\}}.
\end{equation}
With no selected slots, setting $t=1$ still gives the source exponent
$1-2\delta/3$, without any supply assumption.

\paragraph{Scope of the extension.}
The argument proves the changed moment and selected-row interfaces from
the named source inputs.  The source witness and inverse-amplification
statements contain no $7/8$ lower threshold; their relevant restrictions
are respectively $a>51/100$ and the sixth-power-free row condition.
The only zero-free assertion used in the changed fourth moment is
$\beta_*\le(1+\kappa)/2$, explicitly stated above.
The analytic continuation regions for the physical compensated probe,
the changed geometry, and the global endpoint inequality are separate
interfaces and must be verified elsewhere before a new half-plane
theorem is concluded.  This section makes no unsupported extension of
those interfaces.



\section{Endpoint optimization and the bootstrap}
\label{sec:algebra-audit}

This section is an algebraic implication from the row-count and exponent
interfaces stated in the main text.  In particular, its use for zero-free
regions requires those interfaces at every input parameter used below.
Put
\[
 \alpha=\frac56,\qquad
 I_\kappa=\left[\frac{7499}{10000},\frac34\right].
\]
Every terminating decimal used below is an exact rational number.

\subsection{The generalized row envelope}

For \(\kappa\in I_\kappa\), \(0\le x\le1/2\), define
\[
 D_\kappa(x)=3-\left(1+\frac{2}{3\kappa}\right)x,
 \qquad
 P_\kappa(x)=\left(2-\frac{2x}{3\kappa}\right)(1-x).
\]
The selected inverse and plain count exponents are, respectively,
\[
 A_I(r)=1-\delta\{x+(1-x)r\},\qquad
 S_t(r)=1-\delta\left\{\frac{x}{3\kappa}
       +\left(2-\frac{2x}{3\kappa}\right)(t-r)\right\}.
\]
Their intersection gives
\[
 R_{\rm short}(t)=1-\delta+
       \frac{\delta P_\kappa(x)}{D_\kappa(x)}(3/2-t),
 \qquad
 L(t)=1-\delta+(\alpha-\delta)(t-1).
\]
For \(0<\delta\le\kappa\), let
\[
 \mathcal J=(\alpha-\delta)D_\kappa(x)+\delta P_\kappa(x),
 \quad t_*=1+\frac{\delta P_\kappa(x)}{2\mathcal J},
 \quad
 v=\frac{(\alpha-\delta)\delta P_\kappa(x)}{2\mathcal J}.
\]
Then \(1<t_*<3/2\), and exact substitution gives
\[
 R_{\rm short}(t_*)=L(t_*)=1-\delta+v.
\]
This is the minimum of the maximum of these two affine count bounds.

The supply margins also persist on \(I_\kappa\).  Writing
\(c=1/(3\kappa)\), the crossing of \(A_I\) and \(S_t\) occurs at
\[
 r_*(t)=\frac{(2-2cx)t-(1-c)x}{3-(1+2c)x},\qquad
 t-r_*(t)=\frac{(1-x)t+(1-c)x}{3-(1+2c)x}.
\]
For \(1\le t\le3/2\), one has
\[
 r_*(t)\ge r_*(1)\ge\frac{3-c}{5-2c}\ge\frac{23}{37},
 \qquad \frac13\le t-r_*(t)\le\frac12.
\]
Here \(r_*(1)\) decreases with \(x\), whereas
\((3-c)/(5-2c)\) increases with \(c\ge4/9\).
The analogous assertion for \(t-r_*(t)\) follows by differentiating
\((1-cx)/(3-(1+2c)x)\) with respect to \(x\).
Consequently the inverse capacity is at most \(7/37\), its second
strict width margin is at least \(9/37\), and the plain capacity is at
most \(1/(18\kappa)<7/37\).

\subsection{The optimized geometry}

For each \(\kappa\in I_\kappa\), let \(d_\kappa\) be the smaller root of
\begin{equation}\label{eq:audit-root-polynomial}
 p_\kappa(d):=(216\kappa-18)d^2-(468\kappa-57)d+150\kappa-20=0.
\end{equation}
Define
\[
 v_\kappa(d)=
 \frac{(6\kappa-1)d(5/6-d)}
 {2\{(5/6)(15\kappa-2)+(1-9\kappa)d\}},
\]
and set
\begin{equation}\label{eq:audit-optimal-geometry}
 \ell_\kappa=\frac{5-6d_\kappa}{9+18d_\kappa},\qquad
 b_\kappa=
 \frac{(1+3\ell_\kappa)v_\kappa'(d_\kappa)}
      {1-v_\kappa'(d_\kappa)}.
\end{equation}
The remaining geometry is
\[
 l_x=\frac{1-\ell-b}{2},\quad
 l_y=\frac{1-\ell+b}{2},\quad
 h=\frac{1+3\ell+b}{2},\quad
 C=-\frac14+\frac54\ell+\frac b6,
\]
where initially \((\ell,b)=(\ell_\kappa,b_\kappa)\).
With \(y=1/2-x\), the central endpoint exponent is
\begin{equation}\label{eq:audit-endpoint}
 E(\delta,y)=h v+C-\delta(\ell y+b/2).
\end{equation}
At \(y=0\), the quantity \(v\) equals \(v_\kappa(\delta)\).

Here are exact rational enclosures used in the proof.  Put
\(d_-=19429/50000\) and \(d_+=38859/100000\).
Since \(p_\kappa\) is affine in \(\kappa\), the four signs
\[
\begin{array}{c|cc}
 &\kappa=7499/10000&\kappa=3/4\\ \hline
 p_\kappa(d_-)&1875081893/3125000000000&105619/156250000\\
 p_\kappa(d_+)&-15256860787/12500000000000&-715571/625000000
\end{array}
\]
give a root in \((d_-,d_+)\) for every \(\kappa\in I_\kappa\).
The derivative in this rectangle is negative, between \(-183\) and
\(-182\), so this is the smaller root and it is unique there.
In particular,
\begin{equation}\label{eq:audit-parameter-box}
 0.38858<d_\kappa<0.38859,\qquad
 0.1668<\ell_\kappa<0.1669,\qquad
 0.1233<b_\kappa<0.1236.
\end{equation}
For completeness, the intermediate rational interval bounds are
\[
 0.1668348<\ell_\kappa<0.1668405,
 \qquad
 0.075935<v_\kappa'(d_\kappa)<0.076048.
\]
They follow by substituting \(d_-\le d\le d_+\) and
\(\kappa\in I_\kappa\) into
\[
 v_\kappa'(d)=
 \frac{A\{(5/6)B(5/6-2d)+(B-A)d^2\}}
 {2\{(5/6)B+(A-B)d\}^2},\quad
 A=6\kappa-1,\quad B=15\kappa-2.
\]
The last displayed bounds, substituted in
\eqref{eq:audit-optimal-geometry}, give the asserted interval for \(b\).
All denominators in these substitutions are positive.

The optimizing identities are exact.  If
\(W=(5/6)(15\kappa-2)+(1-9\kappa)d\), then
\[
 36W\{v_\kappa(d)-d+1/3\}=p_\kappa(d).
\]
Thus \(v_\kappa(d_\kappa)=d_\kappa-1/3\).  The definitions of
\(\ell_\kappa\) and \(b_\kappa\) now imply
\begin{equation}\label{eq:audit-double-contact}
 E(d_\kappa,0)=0,\qquad
 \partial_\delta E(d_\kappa,0)=0.
\end{equation}

\subsection{The global endpoint inequality}

We prove
\begin{equation}\label{eq:audit-global-endpoint}
 E(\delta,y)\le0
 \quad(0\le\delta\le\kappa,\ 0\le y\le1/2),
\end{equation}
with equality only at \((\delta,y)=(d_\kappa,0)\).
In fact the numerator argument below works for every real \(\delta\).

Put \(c=1/(3\kappa)\) and write
\[
\begin{aligned}
 j_0(y)&=(225-90c)+(90+180c)y,\\
 j_1(y)&=(-162+54c)+(108-216c)y+216cy^2,\\
 P(y)&=1-c/2+2y+2cy^2,\qquad k(y)=b/2+\ell y.
\end{aligned}
\]
Then \(j=j_0+\delta j_1=108\mathcal J\), and multiplication of
\eqref{eq:audit-endpoint} gives the exact identity
\begin{equation}\label{eq:audit-quadratic-numerator}
 -jE=a(y)\delta^2+B(y)\delta+C_0(y),
\end{equation}
where
\[
 a=kj_1+54hP,\qquad
 B=kj_0-Cj_1-45hP,\qquad C_0=-Cj_0.
\]
Write \(a=\sum_{i=0}^3a_i y^i\),
\(B=\sum_{i=0}^2 B_i y^i\), and \(C_0=C_0^{(0)}+C_0^{(1)}y\).
The following table is a rational interval certificate over the entire box
\[
 \ell\in[0.1668,0.1669],\quad b\in[0.1233,0.1236],\quad
 c\in[4/9,10000/22497].
\]
It is obtained directly from the preceding polynomial identities by
interval addition and multiplication, not by sampling parameter values.
\[
\begin{array}{c|rrrr}
 &0&1&2&3\\ \hline
 a_i &[25.56,25.61]&[65.38,65.44]&[46.88,46.93]&[16.01,16.03]\\
 B_i &[-19.92,-19.84]&[-31.51,-31.43]&[-2.14,-2.08]&\\
 C_0^{(i)} &[3.84,3.88]&[3.53,3.57]&&
\end{array}
\]
To make the certificate reproducible without elimination, set
\[
 u_0=225-90c,\ u_1=90+180c,\quad
 w_0=-162+54c,\ w_1=108-216c,\ w_2=216c.
\]
The individual coefficients to which the interval operations apply are
\[
\begin{aligned}
 a_0&=(b/2)w_0+54h(1-c/2),&
 a_1&=(b/2)w_1+\ell w_0+108h,\\
 a_2&=(b/2)w_2+\ell w_1+108hc,&a_3&=\ell w_2,\\
 B_0&=(b/2)u_0-Cw_0-45h(1-c/2),&
 B_1&=(b/2)u_1+\ell u_0-Cw_1-90h,\\
 B_2&=\ell u_1-Cw_2-90hc,&
 C_0^{(0)}&=-Cu_0,\quad C_0^{(1)}=-Cu_1.
\end{aligned}
\]
The constant coefficient of \(4aC_0-B^2\) is zero by
\eqref{eq:audit-double-contact}: the quadratic at \(y=0\) has the
double root \(d_\kappa\).  Its coefficients of \(y,y^2,y^3,y^4\)
are bounded below by, respectively,
\[
 109.7856,\qquad565.1047,\qquad772.9964,\qquad221.4816.
\]
For example, the first lower bound is exactly
\[
 4(25.56\cdot3.53+65.38\cdot3.84)
       -2\cdot19.92\cdot31.51=109.7856.
\]
The remaining three follow from
\[
\begin{aligned}
 [y^2](4aC_0-B^2)&=4(a_1C_0^{(1)}+a_2C_0^{(0)})-B_1^2-2B_0B_2,\\
 [y^3](4aC_0-B^2)&=4(a_2C_0^{(1)}+a_3C_0^{(0)})-2B_1B_2,\\
 [y^4](4aC_0-B^2)&=4a_3C_0^{(1)}-B_2^2.
\end{aligned}
\]
Thus for \(y>0\), \(a(y)>0\) and \(4a(y)C_0(y)-B(y)^2>0\).
Completion of the square proves that the right side of
\eqref{eq:audit-quadratic-numerator} is strictly positive.  At \(y=0\)
it equals \(a_0(\delta-d_\kappa)^2\).
Finally \(j>0\) on the relevant domain: in the \(y\)-coordinates,
\(D=(5/2-c)+(1+2c)y>2\), \(P>3/4\), \(D-P\ge1\), and
\[
 \frac{35}{48}<\mathcal J\le\frac52
 \qquad(0\le\delta\le\kappa\le3/4).
\]
This proves \eqref{eq:audit-global-endpoint} and its equality statement.

\subsection{Strict approximations and the remaining algebraic margins}

Hold \(\kappa,b=b_\kappa\) fixed and replace
\(\ell_\kappa\) by \(\ell_\kappa-\varepsilon\), with
\(0<\varepsilon\le10^{-6}\).  Since \(v\ge0\),
\[
 \frac{\partial E}{\partial\ell}
   =\frac32v+\frac54-\delta y\ge\frac56
 \quad(0\le\delta\le5/6,\ 0\le y\le1/2).
\]
Consequently the perturbed geometry has the uniform strict margin
\begin{equation}\label{eq:audit-strict-margin}
 E\le-\frac56\varepsilon.
\end{equation}
Its output exponent is
\[
 \sigma_0=\frac{11}{12}-\frac\ell4
 =1-\frac2{9+18d_\kappa}+\frac\varepsilon4.
\]
The entire perturbed family remains in the parameter box above.
That box also certifies
\[
\begin{aligned}
 &0.35475\le l_x\le0.35495,
 &&0.4782\le l_y\le0.4784,
 &&0.81185\le h\le0.81215,\\
 &\frac\ell h-\frac7{37}\ge\frac{263}{16243}>0,
 &&l_y-\ell-\frac{11b}{6}\ge\frac{1697}{20000}>0.
\end{aligned}
\]
For the floor bin \(\delta_0=1/50\), the ideal bound \(R=1\)
and \(q\le\delta_0/2\) give
\[
 E_{\rm floor}\le
 C+\frac{\delta_0}{2}(1+3\ell)
 \le-\frac{721}{125000}<0.
\]
The small-row inequality specified in the exponent interface obeys
\[
 h\left(\frac{17}{50}-\frac16\right)-\frac{l_y}{2}+\frac1{50}
 \le-\frac{117491}{1500000}<0.
\]
These are uniform strict algebraic margins.  They do not replace the
analytic hypotheses behind the corresponding source estimates.

\subsection{The bootstrap map and its fixed point}

For an input boundary
\(\sigma\in I_\sigma=[17499/20000,7/8]\), put
\(\kappa=2\sigma-1\) and define
\begin{equation}\label{eq:audit-bootstrap-map}
 f(\sigma)=1-\frac2{9+18d_{2\sigma-1}}.
\end{equation}
Assuming the analytic interfaces, the strict geometries above prove every
output boundary \(f(\sigma)+\varepsilon/4\), and hence, by taking their
intersection, the boundary \(f(\sigma)\) itself.
This limiting step requires no estimate uniform as \(\varepsilon\to0\):
any fixed putative zero to the right is excluded by one sufficiently
small positive \(\varepsilon\).

The root enclosures give
\[
 \frac{349861}{399861}<f(\sigma)<\frac{699731}{799731},
\]
so \(f\) maps \(I_\sigma\) strictly into itself.  Implicit
differentiation of \eqref{eq:audit-root-polynomial} gives
\[
 d_\kappa'=
 -\frac{216d_\kappa^2-468d_\kappa+150}
 {2(216\kappa-18)d_\kappa-(468\kappa-57)}.
\]
On the root rectangle the numerator is between \(3/4\) and \(4/5\),
and the denominator is between \(-183\) and \(-182\).
Thus \(0<d_\kappa'<1/200\), and
\[
 0<f'(\sigma)=\frac{72d_{2\sigma-1}'}{(9+18d_{2\sigma-1})^2}
 <\frac1{500}.
\]
It follows that \(f\) has a unique fixed point \(\sigma_*\) in
\(I_\sigma\).  The iterates beginning at \(\sigma_0=7/8\) decrease
strictly to \(\sigma_*\).  Indeed \(f(\sigma)-\sigma\) is strictly
decreasing and changes sign in the interval, and monotonicity of \(f\)
keeps every iterate above its fixed point.
In particular,
\[
 f(\sigma)<\sigma\quad\hbox{if and only if}\quad
 \sigma\in(\sigma_*,7/8].
\]
Thus, if the analytic argument is formulated using a putative boundary
\(\beta_*\in(\sigma_*,7/8]\), one may apply the geometry directly at
\(\kappa=2\beta_*-1\) and choose
\(0<\varepsilon<4\{\beta_*-f(\beta_*)\}\).
This gives a strictly smaller output boundary and a strictly negative
endpoint margin in one step; an infinite analytic iteration is not
needed for that formulation.

At a fixed point,
\[
 \kappa=2\sigma_*-1=\frac{18d+5}{18d+9}.
\]
Substitution into \eqref{eq:audit-root-polynomial} yields
\[
 p_{(18d+5)/(18d+9)}(d)
 =\frac{3(1188d^3-2160d^2+171d+190)}{18d+9}.
\]
Consequently the exact limiting boundary is
\begin{equation}\label{eq:audit-exact-fixed-point}
 \boxed{\displaystyle
 \sigma_*=1-\frac2{9+18d_*},\qquad
 1188d_*^3-2160d_*^2+171d_*+190=0,
 \quad0.38858<d_*<0.38859.}
\end{equation}
The cubic is strictly decreasing on that interval.  Exact substitution
at the rational endpoints
\(0.38858335426603\) and \(0.38858335426604\) gives opposite signs,
and hence
\[
 \sigma_*\approx0.87495701942009895.
\]
Since \(d\mapsto1-2/(9+18d)\) is increasing and
\[
 \frac{162861}{186136}=1-\frac2{9+18\cdot\frac{162797}{418950}},
 \qquad
 \frac{162797}{418950}>0.38858335426604>d_*,
\]
the rational bound \eqref{eq:main-rational} follows; its margin is
about \(1.25\cdot10^{-9}\).
If each analytic bootstrap implication holds, every zero lies to the
left of every iterate and therefore to the left of \(\sigma_*\).

\subsection{Optimality within the given exponent bounds}

The following final observation is conditional on the broader
low-frequency penalty formula in the main text.  It is an optimality
statement about the existing bounds, not an obstruction to improved
arithmetic estimates.

Allow a nonnegative geometry deficit \(\omega=1-M-\ell\), and put
\[
 p=\frac{(5\ell+2\omega-1)_+}{8},\qquad
 \sigma_0=\frac{11}{12}-\frac\ell4-\frac\omega{12}+p,
 \qquad h=\frac{1+3\ell+b+\omega}{2}.
\]
Suppose the high-frequency exponent is
\[
 E=-\frac14+\frac54\ell+\frac34\omega+\frac b6-p
       -\delta(\ell y+b/2)+hv.
\]
At the same critical point \((\delta,y)=(d_\kappa,0)\), the identity
\(v=d_\kappa-1/3\) makes the coefficient of \(b\) vanish exactly.
The exponent there is
\[
 E=\left(\frac34+\frac32d_\kappa\right)(\ell-\ell_\kappa)
   +\left(\frac7{12}+\frac12d_\kappa\right)\omega-p.
\]
Therefore the necessary condition \(E\le0\) implies
\[
 \boxed{\displaystyle
 \sigma_0-f(\sigma)\ge
 \frac{\omega}{9+18d_\kappa}
 +p\left(1-\frac1{3+6d_\kappa}\right)\ge0.}
\]
Thus allowing \(M+\ell<1\), changing \(b\), or activating this
low-frequency penalty cannot beat the optimized map while these
exponent bounds remain unchanged.  A better boundary requires a new
estimate or an observable with a different exponent ledger.



\section{Analytic estimates for the perturbed probe}
\label{sec:analytic-audit}
\label{app:analytic-audit}

Throughout this section, \(F=\mathbb Q(\sqrt{-3})\), and \(\beta_*\)
denotes the supremum in \texttt{old-eq:1.1b} of the source manuscript:
the supremum of the real parts of nontrivial zeros of \emph{all primitive
finite-order Hecke \(L\)-functions over \(F\)}, with \(1/2\) adjoined.
Poles are excluded.  The supremum must not be replaced by the supremum
for \(\zeta\) alone: the reflected numerator and denominator characters
belong to the larger Hecke family.

We use the exact finite probe of
\texttt{eq:compensated-probe-definition}, with the same slot operation,
zero extensions, excluded primes, and coefficient conventions, but with
variable real geometry
\[
 l_y=l_x+b,\qquad M=2l_x+b,\qquad
 h=1-l_x+\ell,\qquad
 C(s)=s+\frac{l_x}{2}-1+\frac h6.
\]
The slot lengths are positive and sum to \(\ell\).
Define
\[
 \sigma_0=1-\frac h6+\frac b{12},\qquad
 C(\sigma_0)=\frac{l_x}{2}+\frac b{12}.
\]
Here \(M\) is a real log-length, not an inverse polynomial.

\subsection{Explicit sufficient geometric conditions}

The analytic arguments below require
\begin{equation}
\begin{gathered}
 b>0,\quad l_x>\ell>0,\quad h>\tfrac12,\quad
 M+\ell\le1,\quad 1+3\ell-2M\le0,\\
 l_y-\ell-\frac{11b}{6}>0,\qquad \sigma_0>\frac{87}{100}.
\end{gathered}
\label{eq:audit-geometric-conditions}
\end{equation}
In addition, the central arithmetic estimate needs a fixed prime-supply
margin
\begin{equation}
 \frac{\ell}{h+\zeta}>\frac7{37}
 \label{eq:audit-supply}
\end{equation}
for a sufficiently small positive \(\zeta\).  This condition is used only
for \(1/2\le d\le h+\zeta\).  Shorter row ranges use the unselected
witness estimate.

For clarity, the arithmetic inputs to this section are separate from
these geometric conditions.  They are the generalized fourth-moment
lemma with its stated zero-free input, the corresponding selected-witness
row count, and the endpoint inequality described below.  The source
states its fourth-moment lemma only for \(\kappa\ge3/4\); using
\(\kappa<3/4\) requires the extension proved in Section~\ref{sec:extended-moment-audit}.
The present section does not assume that extension merely by changing
a parameter in the source statement.

\subsection{The low estimate with variable geometry}

Fix a rescaled subset of slots of total length \(d\), and put
\[
 0\le d\le\ell,\qquad M'=M-2d,\qquad \ell'=\ell-d.
\]
The proof of \texttt{lem:probe-row-norm}, before its last substitution
of the numerical geometry, gives
\begin{equation}
 \sum_{0<q_m\ll Z^{M'}}|B^J_{m,\sigma}(Z)|^2
 \ll Z^{M'+(1+3\ell-2M+d)_+/4+\epsilon}.
 \label{eq:audit-general-row-energy}
\end{equation}
Here is why its earlier use of the geometry remains valid.  In the
notation of that proof, \(H\le M'-O+\epsilon_Z\) and
\(2A_0\le O+\epsilon_Z\).  The dual-length calculation becomes
\[
 T_d-H\le M'+\ell'-1+O(\epsilon_Z+|\theta_N|)
 =M+\ell-1-3d+O(\epsilon_Z+|\theta_N|)
 \le-3d+O(\epsilon_Z+|\theta_N|).
\]
It follows, exactly as in the source, that the maximum involving the
residual row length is \(H+O(\epsilon)\).  The two subsequent branches
of \texttt{old-eq:5.3} bound the energy by
\[
 M'-\Delta_H+O(\epsilon),\qquad
 M'+\frac{1+3\ell'-2M'-2\Delta_H}{4}+O(\epsilon).
\]
Since \(\Delta_H\ge-O(\epsilon)\), these yield
\eqref{eq:audit-general-row-energy}.  No assumption \(M'\ge1/2\)
enters this calculation.  Positivity of \(M'\), which is all the later
scale argument needs, follows from \(l_x>\ell\) and \(b>0\).

For this subset, the source's additive Gram estimate
\texttt{prop:probe-gram} has
\[
 P_a=\frac{Y'^2}{Q}=q_{b_*}^{-1}Z^b.
\]
Thus \(P_a\ge1\) eventually.  The strict inequality in
\eqref{eq:audit-geometric-conditions} involving \(l_y\) gives
\(Y'/P_a^{11/6}\gg Z^{l_y-\ell-11b/6}\), uniformly in the fixed
rescaled tuple.  Consequently the nonexceptional Gram term
\(P_a^2/Y'\) is bounded by a constant times \(P_a^{1/6}\).
Combining the Gram bound with \eqref{eq:audit-general-row-energy}, then
summing the rescaled prime tuples as in \texttt{old-eq:5.5}, gives
the additional exponent
\[
 -d+\frac{(1+3\ell-2M+d)_+}{8}\le0.
\]
Indeed, if the positive part is nonzero this equals
\((1+3\ell-2M-7d)/8\le0\); otherwise it is \(-d\).
It follows that the same independently defined physical probe satisfies
\begin{equation}
 |I_{\eta,\mathrm{modified}}(Z)|
 \ll_{\eta,\epsilon}Z^{l_x/2+b/12+\epsilon}
 =Z^{C(\sigma_0)+\epsilon}.
 \label{eq:audit-low}
\end{equation}
More generally, without the condition \(1+3\ell-2M\le0\), but keeping
\(M+\ell\le1\) and the other scale and Gram conditions, the same proof
gives the additional low exponent
\[
 p=\frac{(1+3\ell-2M)_+}{8}.
\]
Indeed \(-d+(1+3\ell-2M+d)_+/8\) is nonincreasing in \(d\): its
two possible slopes are \(-1\) and \(-7/8\).  Its maximum at \(d=0\)
is exactly \(p\).  If \(\mathfrak d=1-M-\ell\ge0\), the resulting
low threshold and central endpoint are consequently
\[
 \sigma_0=\frac{11}{12}-\frac\ell4-\frac{\mathfrak d}{12}+p,
 \qquad p=\frac{(5\ell+2\mathfrak d-1)_+}{8},
\]
\[
 E(h)=-\frac14+\frac{5\ell}{4}+\frac{3\mathfrak d}{4}
       +\frac b6-p-\delta(\ell y+b/2)+h v,
 \qquad R=1-\delta+v.
\]
These formulas justify the general low-penalty ledger used in the
geometry-optimality argument; they do not assert optimality among
different arithmetic moment estimates.

\subsection{Extension of the principal Euler region}

The displayed local formulas in \texttt{lem:local-euler} extend its
second region to
\begin{equation}
 \Re s\ge\frac{87}{100},\qquad
 \Re w\ge\frac{19}{20},\qquad
 \Re z\ge\frac{33}{200}.
 \label{eq:audit-principal-region}
\end{equation}
This is a direct estimate, not a continuity assertion.  The source's
local quantities satisfy
\[
 |R|\le Q^{-221/100},\qquad |V|\le Q^{-99/100},\qquad
 |D|\le Q^{-87/100}.
\]
Using \texttt{eq:local-H-defect} and the displayed table of \(J_j\),
the nearest unramified defect is
\(Q^{1-\Re s-\Re w-6\Re z}\le Q^{-181/100}\), and the nearest
ramified defect is \(Q^{1-\Re s-\Re w}\le Q^{-41/50}\).
The other unramified terms are bounded, for example, by
\(Q^{-182/100}\), \(Q^{-186/100}\), or \(Q^{-221/100}\).
The other ramified boundary terms have exponents at most
\[
 -111/100,\quad -148/100,\quad -156/100,\quad -193/100,\quad -222/100.
\]
Therefore, uniformly in all heights and unit phases,
\begin{equation}
 H_p-1=O(Q^{-181/100})\ (p\nmid u),\qquad
 H_p-1=O(Q^{-41/50})\ (p\mid u).
 \label{eq:audit-euler-defects}
\end{equation}
Normal convergence and the divisor-product argument of the source give
\(\mathcal H_{\eta,u}\ll_\epsilon q_u^\epsilon\) in
\eqref{eq:audit-principal-region}.  For \(u=1\), deleting primes up to
\(P_0\) gives \(\mathcal H_{\eta,1}=1+O(P_0^{-4/5})\): the stronger
tail exponent supplied by \eqref{eq:audit-euler-defects} is \(81/100\).

The source's quotient-free formula \texttt{old-eq:5.13c} for \(G_p\)
continues the full selected-tuple correction into this region.  There
is no division by a possibly vanishing \(H_p\) or \(P_p\) in that
formula.  For the principal row, after taking \(P_0\) large enough that
\(|H_p-1|<1/2\), \texttt{old-eq:5.13b} gives
\[
 G_p/H_p=-1+O(Q^{-87/100}).
\]
The contributing error powers are bounded by
\(Q^{-87/100},Q^{-99/100},Q^{-134/100},Q^{-94/100}\).
The fixed-ray prime normalizer \texttt{lem:ray-prime-normalizer}
therefore proves, for the source's \(A_T(Z)\),
\begin{equation}
 |A_T(Z)|\asymp_\eta(\log Z)^{-K},\qquad
 \mathfrak H_{\eta,1,Z}(s,1,1/6)
 =H_\eta(s)Z^{\ell/6}A_T(Z)(1+O(Z^{-\mu})),
 \label{eq:audit-normalizer}
\end{equation}
uniformly on \(\Re s=\beta_*+e\), for any fixed
\(0<\mu<(87/100)\min_i\ell_i\).  Here
\(H_\eta(s)=\mathcal H_{\eta,1}(s,1,1/6)\).
Increasing the same fixed \(P_0\) ensures
\[
 \sup_{\Re s>\sigma_0}|H_\eta(s)-1|\le\tfrac12.
\]
All excluded sets and lower thresholds may depend on the fixed target;
the exponent \(\mu\) is chosen before it.  Only a fixed-ray prime
asymptotic is used, not a uniform moving-conductor prime theorem.

\subsection{The other contours and the exact central exponent}

The exact Poisson identity \texttt{eq:compensated-poisson-identity}
still holds for this physical probe, because at each fixed \(Z\) the
modification is a finite linear combination of the same absolutely
convergent rescaled base identities.  Its weight has the general form
in \texttt{eq:stage-mellin-exponent}.  In particular it requires
\(h=1-l_x+\ell\), not the original numerical values of the geometry.

The first Euler region \texttt{eq:local-region-one} and the dynamic
local-error estimate \texttt{prop:probe-errors} are unchanged.  The
proof of \texttt{lem:fixed-bin-contour} uses only this first region,
the buffered zero rectangles, and the global line to the right of
\(\beta_*\).  Its stated assumption \(\sigma_0\ge7/8\) is not used
elsewhere in that proof.  The same is true of the algebraic accounting
in \texttt{lem:bin-contour-accounting}.

For a bin \(a=(1+\delta)/2\), prime-amplitude mean \(q\), and row-count
exponent \(R\), that accounting gives exactly
\begin{equation}
 E(d)=a-\sigma_0+h(17/50-1/6)-a l_y-\ell/2+q\ell
          +d(R+\delta/2-17/50).
 \label{eq:audit-central-exponent}
\end{equation}
At \(d=h\), writing \(y=1/2-q/\delta\), it becomes
\begin{equation}
 E(h)=h(R+\delta+1/2)-1-\ell(1+\delta y)
                  -b(7/12+\delta/2).
 \label{eq:analytic-endpoint}
\end{equation}
The generalized arithmetic endpoint lemma is required to prove
\(E(h)\le0\), uniformly for the relevant \(\delta,y\).  It is sufficient
that this inequality be non-strict.  Under the contradiction
\(\Delta=\beta_*-\sigma_0>0\), it gives a strict saving \(\Delta\)
relative to \(C(\beta_*)\).  The extension to \(d\le h+\zeta\), the
slot rounding, the detector padding, and all amplitude losses can be
chosen to consume less than \(\Delta/2\).  For \(d\le h\), the
frequency slope is positive whenever \(R\ge1-\delta\) and
\(\delta\le3/4\), since
\[
 R+\delta/2-17/50\ge33/50-\delta/2\ge57/200.
\]

The floor \(\delta_0=1/50\) uses \(R=1\), without a witness.  Its
endpoint is bounded above by
\begin{equation}
 E_{\rm floor}(h)
 =h(3/2+1/50)-1-\ell-b(7/12+1/100).
 \label{eq:audit-floor}
\end{equation}
For \(d\le1/2\), the unselected estimate
\(R=76/75-2\delta/3\) is available, including the floor.  Using
\(q\le\delta/2\), its worst endpoint is the larger value, for
\(\delta\in[1/50,3/4]\), of the affine function
\begin{equation}
 \frac12-\sigma_0+h(17/50-1/6)-\frac{l_y+\ell}{2}
 +\frac12(76/75-17/50)
 +\delta\left(\frac5{12}+\frac{\ell-l_y}{2}\right).
 \label{eq:analytic-short}
\end{equation}
Thus checking its two endpoints is sufficient for the whole short range.

The principal-residue proof \texttt{lem:principal-residue-interface}
now takes place in \eqref{eq:audit-principal-region}; its two scalar
poles and their residues are unchanged.  By \eqref{eq:audit-normalizer},
the three resulting errors, after normalization, have exponents
\[
 C(\beta_*)+(1+h)e-l_y/20+\epsilon,\qquad
 C(\beta_*)+e-h/600+\epsilon,\qquad
 C(\beta_*)+e-\mu+\epsilon.
\]
These are strict savings once the positive padding is sufficiently small.

For small nonprincipal rows the source's
\texttt{lem:compensated-absolute-local-bounds} remains valid with
\(\Re s\ge87/100\).  Off \(u\), its four error exponents are at most
\[
 -87/100,\quad -51/25,\quad -239/100,\quad -77/50;
\]
on \(u\), the largest term is
still \(Q^{1/2}\).  Hence the positive tuple sum is
\(O_\epsilon(U^\epsilon Z^{17\ell/50})\).
In the source proof of \texttt{lem:outer-row-tails}, replace the claimed
path in \(\mathcal D_1(3/8)\) by one in \(\mathcal D_1(1/3)\).
This is legitimate because \(\beta_*>\sigma_0>87/100>5/6\).
For \(d_{\min}=1/100\), the resulting small-row exponent relative to
\(C(\beta_*)\) is
\begin{equation}
 h(17/50-1/6)-l_y/2+\frac{63}{5000}+e+\epsilon.
 \label{eq:audit-small}
\end{equation}
For large rows, the source's absolute lines \((s,w,z)=(2,2,v)\),
\(v>2\), are unchanged.  With \(B_0=l_x/2+1+l_y\), their sum for
\(U>Z^{h+\zeta}\) is bounded by
\[
 Z^{B_0+(h+\zeta)(1+\epsilon)-\zeta v+\epsilon}.
\]
A fixed sufficiently large \(v\) makes it smaller than any required
power.  The order is: choose \(\zeta\), then \(v\), both before the
target and the eventual scale.

\subsection{Verification on the proposed geometry interval}

The optimized family in the main proof obeys the rational enclosures
\[
 \frac{16683}{100000}<\ell<\frac{16685}{100000},\qquad
 \frac{123}{1000}<b<\frac{124}{1000},\qquad
 M=1-\ell,\qquad h=\frac{1+3\ell+b}{2}.
\]
They imply
\[
 0.811745<h<0.812275,\quad
 0.354575<l_x<0.355085,\quad
 0.478075<l_y<0.478585,
\]
and \(\sigma_0=11/12-\ell/4>0.8749>87/100\).
In particular \(1+3\ell-2M=5\ell-1<-0.1657\),
\(l_x-\ell>0.1877\),
\(l_y-\ell-11b/6>0.0843\), and
\(\ell/h>0.205>7/37\).  Thus
\eqref{eq:audit-geometric-conditions} and
\eqref{eq:audit-supply} hold with uniform positive margins.

For this family \eqref{eq:audit-floor} simplifies to
\[
 -\frac6{25}+\frac{32\ell}{25}+\frac b6<-0.0057.
\]
The coefficient of \(\delta\) in \eqref{eq:analytic-short} is positive.
At \(\delta=3/4\), that expression simplifies to
\[
 -\frac{871}{1200}+\frac{243h}{200}
                  -\ell-\frac{23b}{24}
 =-\frac{71}{600}+\frac{329\ell}{400}-\frac{421b}{1200}
 <-0.0242.
\]
Finally \eqref{eq:audit-small}, before its adjustable losses, is
less than \(-0.0856\).  The principal geometric margins satisfy
\(l_y/20>0.0239\) and \(h/600>0.00135\).
These inequalities verify every outer and floor interface without an
appeal to qualitative robustness.  The nonfloor endpoint remains the
explicit arithmetic polynomial inequality proved separately.

\subsection{Mellin continuation and order of choices}

Set \(J_\eta=I_{\eta,\mathrm{modified}}/(c_S A_T)\), with the same
normalizer in the low and high estimates.  The positive scalar \(c_S\)
is the product of the two principal residues from the source.  Equations
\eqref{eq:audit-low} and \eqref{eq:audit-normalizer} give
\(J_\eta\ll Z^{C(\sigma_0)+\omega}\) for any fixed
\(0<\omega<\Delta\).  The central endpoint estimate and the explicit
margins above give a target-independent \(m>0\) and
\[
 |J_\eta-f_\eta|
 \ll_{\eta,N}Z^{C(\beta_*)-m}(1+T_1)^{A_\eta}
                  +Z^{B_\eta}T_1^{-N},
\]
where \(A_\eta,B_\eta\) are finite and fixed before the external tail
order \(N\).  This is exactly the contract in
\texttt{lem:late-height-closure}.  Its proof chooses
\(T_1=Z^{\tau_\eta}\) after \(\eta\), with
\(\tau_\eta\le m/[4(A_\eta+1)]\) and the detector's additional ceiling,
and then chooses \(N\).  It supplies the common saving \(m/2\).
The physical probe and its principal signal do not depend on \(T_1\).

For completeness, the signal is
\[
 f_\eta(Z)=\frac1{2\pi i}\int_{(2)}
 Z^{C(s)}e^{(s-5/6)^2}\frac{H_\eta(s)}{L_F^S(s,\eta)}\,ds.
\]
The common low and high savings imply
\(f_\eta(Z)\ll Z^{C(\beta_*)-\varepsilon_*}\), with
\(0<\varepsilon_*<\Delta\) independent of \(\eta\), after shrinking
the saving if necessary. Thus the half-plane to be continued remains
inside \(\Re s>\sigma_0\), where the multiplier is nonvanishing.
Moving only the
absolutely convergent line \(\Re s=2\) to arbitrary larger real parts
gives arbitrarily rapid power decay as \(Z\downarrow0\).
Consequently
\[
 \int_0^\infty f_\eta(Z)Z^{-C(s)}\,\frac{dZ}{Z}
\]
is holomorphic for \(\Re s>\beta_*-\varepsilon_*\).  Fourier inversion
on \(\Re s=2\), followed by the identity theorem, identifies it with
\(e^{(s-5/6)^2}H_\eta(s)/L_F^S(s,\eta)\) where \(\Re s>1\).
Since \(|H_\eta|\ge1/2\), it continues \(1/L_F^S\) holomorphically
to the larger half-plane.  This contradicts the definition of
\(\beta_*\), exactly as in \texttt{lem:continuation-criterion}.
The deleted factors have no zeros for \(\Re s>0\).

Thus the analytic interfaces prove the claimed new half-plane once the
generalized arithmetic moment, row-count, and nonfloor endpoint lemmas
have been established.  The transfer to all finite-order Hecke
characters and to Dirichlet \(L\)-functions is the source's
\texttt{prop:hecke-dirichlet-transfer}; principal poles at one are allowed.
For a limiting bootstrap value, apply this implication at each finite
stage and take the infimum of the resulting bounds on the same
\(\beta_*\).  One must not claim a fresh strict endpoint saving at the
limiting tangency without the contradiction parameter \(\Delta>0\).

\section*{AI assistance}
The author found this somewhat accidentally. He asked GPT-6 Astra to try
to apply the results of OpenAI's $7/8$ proof to some independent
investigations of the Riemann hypothesis that he was pursuing. It did
the opposite, and applied some optimization ideas he had tried to the
OpenAI proof. Somehow, this worked and produced a very slight
improvement.

\begin{thebibliography}{9}
\bibitem{OA}
OpenAI, \emph{The Quasi-Riemann Hypothesis: A Zero-Free Half-Plane
$\Rea s>7/8$}, 30 September 2026. Version of 6 October 2026,
\href{https://github.com/openai/math/blob/adc7f1241b42e322a6451854ab7e4b4c146bf78a/preprints/The-Quasi-Riemann-Hypothesis-September-30-2026/build/paper.tex}{source TeX at GitHub (commit \texttt{adc7f124})}.
\end{thebibliography}

\end{document}
